应用不是在理论末尾贴几个行业名词。一个好的案例应让变量、损失、约束、正则化、算法和验证重新汇合，并暴露理论假设与现实数据之间的距离。

本单元选择四个彼此呼应的案例：

\begin{enumerate}
\tightlist
\item
  光滑经验风险与 \(L_2\) 正则化比较 Ridge 与 logistic 回归，说明损失函数怎样决定曲率、概率解释与数值实现。
\item
  \(L_1\) 正则为何产生稀疏通过 Lasso 的几何、KKT 和软阈值解释精确零从何而来，也讨论相关特征下的局限。
\item
  最大间隔与支持向量机从分离超平面出发推导硬间隔、hinge loss 和对偶支持向量。
\item
  风险、收益与投资组合把二次规划、对偶价格和估计误差放进同一个决策问题。
\end{enumerate}

这些模型都是凸的，但``全局最优''只说明已经忠实求解所写下的模型，不保证模型假设、数据分布或评价指标正确。每节最后都会把优化证书与统计、业务验证区分开来。

\subsection{一笔账看清正则化的三重角色}

用 10 个样本拟合 \(y=3x_1+2x_2+\varepsilon\)，其中 \(x_1\) 和 \(x_2\) 的相关系数为 0.98。最小二乘的系数为 \((4.1,-1.0)\)，符号相反且数值放大——共线性让协方差接近奇异。Ridge（\(\lambda=1\)）得 \((2.1,1.8)\)，数值稳定且接近真实 \((3,2)\)。Lasso（\(\lambda=0.5\)）可能因相关特征将其中一个压为 0，得到 \((3.5,0)\)，解并非不正确，但选出的变量不稳定。一次实验中正则化同时影响了条件数、解的范数和变量选择，三重效果不能只归给``防过拟合''一个标签。

\subsection{怎么读这一章}

四篇都从``目标函数的每一项在现实里买来了什么''入手。平方正则控制整体幅度，绝对值正则鼓励系数归零，最大间隔用几何余量换稳健分类，投资组合则把收益与风险放进同一个权衡。算法只是最后一步；先读懂建模选择。所需算法可回看 \hyperref[ch:076]{``凸优化算法：利用梯度、曲率与可分结构''}。

\subsection{本章篇目}

以下篇目按概念依赖排列。

\begin{itemize}
\tightlist
\item \hyperref[sec:09-Convex-Optimization/07-Applications/01]{``\texorpdfstring{光滑经验风险与 \(L_2\) 正则化}{光滑经验风险与 L_2 正则化}''}；
\item \hyperref[sec:09-Convex-Optimization/07-Applications/02]{``\texorpdfstring{\(L_1\) 正则为何产生稀疏}{L_1 正则为何产生稀疏}''}；
\item \hyperref[sec:09-Convex-Optimization/07-Applications/03]{``最大间隔与支持向量机''}；
\item \hyperref[sec:09-Convex-Optimization/07-Applications/04]{``风险、收益与投资组合''}。
\end{itemize}
